% =========================================================
%  部件骨架模板 part_qN.tex（按题拆分的论文部件）
%
%  用法：按题目数量 N 生成 N+1 份（part_q1.tex … part_qN.tex + part_common.tex），
%        每个成员复制一份自己负责的部件，只填写自己负责的片段，
%        然后在组装空间由 AI 拼接进 main.tex。
%
%  红线（违反会导致组装失败）：
%  1. 本文件内禁止 \documentclass / \begin{document} / \usepackage / \newenvironment
%     ——这些只存在于 main.tex。本文件是「片段集合」，不是独立文档。
%  2. label 命名空间：\label{q1:fig:xxx} / \label{q1:tab:xxx} / \label{q1:eq:xxx}
%    （公有部件用 \label{cmn:...}）。禁止跨题文字引用（如"如问题三所示"）。
%  3. 内容必须自包含：不依赖其他部件中定义的文字/内容；\ref 可跨部件引用
%     （label 前缀保证了不冲突），但文字叙述不得引用其他部件的内容。
%
%  预览：用 preview_qN.tex（模板 + \input{part_qN.tex}）独立编译查看效果。
%  配套背景材料：part_qN_background.md（该题原文/硬约束/假设/分析/建模算法说明，
%        分发时随附给队友参考，不进论文、不参与组装；模板见
%        references/assembly/part-background-template.md）
% =========================================================

% ─── Q1-摘要句组（一段话，自包含，不写跨题连接词）───────────
针对问题一，建立了【模型名】模型，采用【算法】求解，得到【量化结果】。

% ─── Q1-问题重述 ─────────────────────────────────────────
% 问题一的背景与任务描述（用自己的话，1~2 段）

% ─── Q1-问题分析 ─────────────────────────────────────────
% 问题一的建模思路分析：关键变量、约束条件、数据特征、方法选择理由

% ─── Q1-模型建立 ─────────────────────────────────────────
% 变量定义 → 目标函数 → 约束条件 → 求解思路
% 公式示例（编号连续，组装后由 LaTeX 统一管理）：
% \begin{equation}
%     \min_{x} \quad f(x) = \sum_{i=1}^{n} w_i g_i(x)
%     \label{q1:eq:obj}
% \end{equation}

% ─── Q1-模型求解 ─────────────────────────────────────────
% 求解算法与步骤（1~2 段 + 枚举）

% ─── Q1-结果分析 ─────────────────────────────────────────
% 结果表格（三线表 + \setlength{\tabcolsep}{} + 表标题在上）：
% {\renewcommand{\arraystretch}{1.38}
% \setlength{\tabcolsep}{6pt}
% \begin{longtable}{ccc}
%     \caption{问题一计算结果}\label{q1:tab:result}\\
%     \toprule[1.5pt]
%     ...  & ... & ... \\
%     \midrule[1pt]
%     \endfirsthead
%     \endhead
%     \bottomrule[1.5pt]
%     \endlastfoot
%     ...  & ... & ... \\
% \end{longtable}}
%
% 结果图（图标题在下，\cref 引用不加"图"字前缀）：
% \begin{figure}[H]
%     \centering
%     \includegraphics[width=0.7\textwidth]{figures/q1_result.pdf}
%     \caption{问题一结果图示}
%     \label{q1:fig:result}
% \end{figure}

% ─── Q1-灵敏度（若属于本题）─────────────────────────────
% 本题专属的灵敏度/稳健性分析（公有灵敏度放 part_common.tex）

% ─── Q1-附录（代码/补充表格，若有）───────────────────────
% \lstinputlisting[frame=tb]{appendix_q1.py}（代码文件随部件一起交付）
% 注：附录章节在组装时统一归入附录部分，表/图编号由 LaTeX 全局管理

% =========================================================
%  片段内容到此结束——本部件文件不得出现任何文档级命令
% =========================================================
